\documentclass[../DoAn.tex]{subfiles}
\begin{document}
\sloppy

Chương này trình bày các công nghệ và kiến trúc được lựa chọn để xây dựng hệ
thống quản lý ký túc xá. Trong quá trình thực hiện đồ án, nhóm đã cân nhắc
nhiều lựa chọn khác nhau và đưa ra quyết định dựa trên yêu cầu cụ thể của bài
toán: cần xử lý các đợt phân bổ phòng theo chu kỳ học kỳ, hỗ trợ đồng thời
cả web lẫn ứng dụng di động, đảm bảo thông báo thời gian thực đến sinh viên,
và duy trì tính nhất quán dữ liệu khi nhiều quản trị viên thao tác cùng lúc.

% =====================================================================
\section{Kiến trúc hệ thống}
\label{sec:architecture}
% =====================================================================

Hệ thống được tổ chức theo kiến trúc phân lớp~\cite{fowler2002patterns}, trong
đó backend chịu toàn bộ trách nhiệm xử lý nghiệp vụ và truy cập cơ sở dữ liệu,
frontend hiển thị theo mô hình server-side rendering, và một nhóm script độc
lập phục vụ các tác vụ hàng loạt như khởi tạo dữ liệu, cấu hình môi trường
và migration.

Lý do chọn kiến trúc này xuất phát từ đặc điểm của nghiệp vụ phân bổ phòng:
mỗi đợt phân bổ cần đi qua nhiều bước theo thứ tự xác định --- kiểm tra điều
kiện, tính điểm ưu tiên, ghi nhận kết quả --- và không thể thực thi song song
trên cùng tập sinh viên. Kiến trúc phân lớp giúp tách rõ phần tính toán điểm
ưu tiên ở tầng Service khỏi phần truy vấn dữ liệu ở tầng Model, qua đó dễ
kiểm soát và kiểm thử từng phần độc lập.

Luồng xử lý chính hoạt động như sau: client gửi yêu cầu qua trình duyệt hoặc
ứng dụng di động, server xác thực danh tính và cấp token hoặc session, tầng
Service thực hiện tính điểm ưu tiên và ghi kết quả vào cơ sở dữ liệu, sau đó
phản hồi về client. Thông báo thời gian thực được đẩy đến sinh viên qua kênh
Socket.IO~\cite{socketio2024docs} ngay sau khi kết quả phân bổ được ghi nhận.
Các tác vụ chạy lâu như seed dữ liệu hoặc migration được tách hoàn toàn ra
script riêng, không ảnh hưởng đến hiệu năng API chính.

Một điểm quan trọng trong thiết kế là hệ thống lưu lại snapshot chính sách
phân bổ tại thời điểm sinh viên nộp đơn. Điều này đảm bảo khi ban quản lý
thay đổi tiêu chí ưu tiên sau đó, các đơn đã nộp vẫn được xét theo chính sách
cũ --- đảm bảo tính công bằng và khả năng truy vết lý do phân bổ.

% =====================================================================
\section{Công nghệ backend}
\label{sec:backend}
% =====================================================================

\subsection*{Runtime và framework}

Backend được xây dựng trên Node.js~\cite{nodejs_docs} phiên bản 22 LTS với
framework Express.js~\cite{expressjs_docs} 4.x. Node.js được chọn vì toàn bộ
hệ thống dùng JavaScript cả frontend lẫn backend, giúp tái sử dụng logic xử
lý dữ liệu và giảm chi phí chuyển đổi ngôn ngữ trong nhóm phát triển. Vòng
lặp sự kiện đơn luồng của Node.js phù hợp với đặc trưng I/O-bound của hệ
thống ký túc xá --- phần lớn tải đến từ truy vấn cơ sở dữ liệu và giao tiếp
mạng, không phải tính toán nặng về CPU. Express.js được chọn vì tính tối giản
--- không áp đặt cấu trúc cứng nhắc, cho phép tổ chức route và middleware theo
cách phù hợp với từng nhóm nghiệp vụ cụ thể như chu kỳ phân bổ, quản lý vi
phạm hay kế thừa chính sách.

\subsection*{Cơ sở dữ liệu}

Cơ sở dữ liệu sử dụng MongoDB Atlas~\cite{mongodb2024datamodel} kết nối qua
Mongoose~\cite{mongoose_docs} 8.x làm ODM (Object--Document Mapper). Mô hình
document của MongoDB phù hợp với dữ liệu sinh viên và phòng ở vốn có cấu trúc
thay đổi theo từng học kỳ --- ví dụ, tiêu chí ưu tiên có thể bổ sung trường
mới mà không cần migration phức tạp như cơ sở dữ liệu quan hệ. Mongoose cung
cấp schema validation tại tầng ứng dụng, các middleware kiểm tra trước khi lưu
(pre-save hooks), và phương thức truy vấn rõ ràng, đảm bảo ràng buộc dữ liệu
chặt chẽ mà không phụ thuộc hoàn toàn vào cơ sở dữ liệu.

Hệ thống định nghĩa hơn 19 schema được tổ chức trong thư mục
\texttt{src/schemas/}, bao gồm các schema nghiệp vụ cốt lõi: chu kỳ phân bổ,
nhật ký kiểm toán, vi phạm nội quy, yêu cầu bảo trì, thông báo, khai báo ưu
tiên, refresh token di động và hàng đợi sự kiện outbox. Mỗi schema được tách
thành file riêng để dễ bảo trì và kiểm thử độc lập.

\subsection*{Ghi log và giám sát}

Hệ thống tích hợp Winston~\cite{winston_docs} làm thư viện logging có cấu trúc,
ghi log dạng JSON kèm metadata đầy đủ cho mọi thao tác quan trọng như phân bổ
phòng, thay đổi chính sách và xử lý vi phạm. Mức log được điều khiển qua biến
môi trường \texttt{LOG\_LEVEL}. Sentry~\cite{sentry_docs} (\texttt{@sentry/node}
10.x) được khởi tạo ngay khi ứng dụng bắt đầu để tự động thu thập error events
và performance traces, hỗ trợ phát hiện sự cố trong môi trường production.
Ngoài ra, thư viện \texttt{@opentelemetry/api}~\cite{opentelemetry_docs} được
tích hợp để cung cấp nền tảng quan sát (observability) theo chuẩn OpenTelemetry,
cho phép gắn span và trace vào các luồng nghiệp vụ quan trọng.

\subsection*{Bảo mật tầng trung gian}

Về bảo mật tầng trung gian, hệ thống dùng Helmet~\cite{helmet_docs} 8.x để
thiết lập các HTTP security header, express-rate-limit~\cite{express_rate_limit}
8.x để giới hạn số yêu cầu nhằm ngăn chặn tấn công brute-force, và
express-mongo-sanitize để chặn NoSQL injection. Dữ liệu đầu vào được kiểm
tra bằng cả express-validator~\cite{express_validator_docs} và
Joi~\cite{joi_docs}, tùy theo ngữ cảnh: Joi dùng cho các cấu hình cần schema
phức tạp, express-validator dùng cho kiểm tra nhanh ở tầng route. Trước khi
ứng dụng khởi động, module \texttt{validateEnv.js} kiểm tra toàn bộ biến môi
trường bắt buộc; nếu thiếu bất kỳ secret nào, ứng dụng thoát ngay với exit
code 1 thay vì chạy trong trạng thái thiếu cấu hình.

\subsection*{Lưu trữ tệp và xuất dữ liệu}

Để lưu trữ ảnh và tài liệu, hệ thống dùng
Cloudinary~\cite{cloudinary_docs} kết hợp Multer~\cite{multer_docs} và
\texttt{multer-storage-cloudinary}. Ảnh ký túc xá, minh chứng kèm yêu cầu dịch
vụ và media cho tính năng xem phòng 360 độ đều được đẩy thẳng lên Cloudinary
mà không cần lưu tạm trên đĩa máy chủ. ExcelJS~\cite{exceljs_docs} được dùng
để xuất kết quả phân bổ và số liệu thống kê ra file \texttt{.xlsx} phục vụ ban
quản lý. Thư viện \texttt{qrcode}~\cite{qrcode_docs} sinh mã QR cho thẻ cư
trú sinh viên phía server.

\subsection*{Cơ chế truyền sự kiện}

Một điểm đáng chú ý là hệ thống thiết kế cơ chế truyền sự kiện (event
transport) có thể hoán đổi được giữa ba lựa chọn: in-process cục bộ (mặc định),
Redis pub/sub~\cite{redis2024pubsub} và Apache Kafka~\cite{kafka_docs}, điều
khiển qua biến môi trường \texttt{EVENT\_TRANSPORT}. Trong cấu hình hiện tại,
hệ thống chạy ở chế độ local (in-process) để đơn giản hóa triển khai; Redis
pub/sub được kích hoạt khi cần đồng bộ sự kiện giữa nhiều instance. Nhờ thiết
kế cắm được này, có thể nâng cấp sang Kafka khi quy mô mở rộng mà không cần
sửa mã nghiệp vụ.

% =====================================================================
\section{Công nghệ frontend}
\label{sec:frontend}
% =====================================================================

Giao diện web --- bao gồm cả admin portal lẫn student portal --- được xây dựng
theo kiến trúc đa trang sử dụng EJS~\cite{ejs_docs} làm template engine phía
server, kết hợp HTML5, CSS3 thuần và Bootstrap~\cite{bootstrap_docs} 5.x.
Biểu tượng giao diện sử dụng Font~Awesome~6.x~\cite{fontawesome_docs} nhúng
qua CDN.

Quyết định dùng EJS thay vì framework hiện đại như React hay Vue xuất phát từ
đặc điểm của hệ thống: phần lớn trang cần render dữ liệu phân bổ và trạng thái
phòng ngay tại server trước khi gửi về trình duyệt. Cách này đơn giản hóa đáng
kể logic phía client, đặc biệt với các biểu mẫu đăng ký phức tạp nhiều bước và
trang quản trị cần hiển thị số lượng lớn bản ghi.

Tính năng bản đồ hiển thị vị trí ký túc xá sử dụng OpenStreetMap~\cite{openstreetmap_docs}
theo phương pháp nhúng iframe trực tiếp từ dịch vụ công cộng
\texttt{openstreetmap.org/export/embed.html}. Phương pháp này không cần thư viện
JavaScript bổ sung, không tạo ra kết nối mạng từ phía client ngoài nội dung
iframe, và hoàn toàn tương thích với chính sách Content Security Policy của hệ
thống. Vị trí marker và khung nhìn được xác định qua tham số URL dựa trên tọa
độ lưu trong cơ sở dữ liệu.

Hệ thống áp dụng nguyên tắc progressive enhancement: các trang cơ bản vẫn hoạt
động khi JavaScript bị tắt, các tương tác nâng cao được bổ sung dần bằng các
module JavaScript nhỏ tự viết, không phụ thuộc framework phía client nào. Toàn
bộ dữ liệu đầu vào đều được kiểm tra phía server trước khi xử lý.

% =====================================================================
\section{Xác thực và bảo mật}
\label{sec:security}
% =====================================================================

Hệ thống áp dụng hai cơ chế xác thực song song tùy theo loại client. Web portal
dùng session-based authentication qua \texttt{express-session}~\cite{express_session_docs},
lưu session ID trong cookie với thuộc tính \texttt{HttpOnly}, \texttt{SameSite=strict}
và thời hạn 24 giờ. Mobile API dùng JWT~\cite{rfc7519} stateless authentication
với hai loại token: access token TTL 15 phút và refresh token TTL 30 ngày, ký
bằng hai secret riêng \texttt{MOBILE\_JWT\_ACCESS\_SECRET} và
\texttt{MOBILE\_JWT\_REFRESH\_SECRET}, quản lý bằng thư viện
\texttt{jsonwebtoken}~\cite{jsonwebtoken_docs}.

Hệ thống dùng Passport.js~\cite{passportjs_docs} 0.7.x làm tầng trừu tượng
hóa xác thực cho web portal. \texttt{LocalStrategy} xử lý đăng nhập truyền
thống bằng định danh và mật khẩu đã băm bằng bcrypt~\cite{bcrypt_npm} 6.x.
\texttt{GoogleStrategy}~\cite{passport_google_docs} và
\texttt{MicrosoftStrategy}~\cite{passport_microsoft_docs} theo chuẩn OAuth
2.0~\cite{rfc6749} cho phép sinh viên đăng nhập bằng tài khoản tổ chức mà
không cần tạo mật khẩu riêng. Việc tách chiến lược xác thực ra khỏi logic
nghiệp vụ giúp dễ bổ sung nhà cung cấp danh tính mới mà không ảnh hưởng đến
phần còn lại.

Mã QR thẻ cư trú dùng HMAC-SHA256 để ký payload chứa thông tin sinh viên và
timestamp, với TTL 24 giờ và rotate mỗi lần gọi API. Ứng dụng di động chỉ
nhận chuỗi token và hiển thị QR, không thực hiện thao tác mật mã nào phía
client. Hệ thống phát hiện bất thường token: khi một refresh token bị dùng lại
từ thiết bị khác hoặc trùng \texttt{jti}, cơ chế risk score được kích hoạt và
có thể tự động vô hiệu hóa toàn bộ phiên. Mọi thao tác trên dữ liệu nhạy cảm
đều được ghi audit log chi tiết.

Xác thực hai yếu tố được hiện thực bằng Speakeasy~\cite{speakeasy_npm} 2.x
theo chuẩn TOTP~\cite{rfc6238}: server sinh khóa bí mật cho từng tài khoản,
mã hóa thành QR code để người dùng đăng ký vào ứng dụng xác thực như Google
Authenticator, sau đó xác minh mã sáu chữ số sinh theo chu kỳ thời gian.
Xác thực hai yếu tố bắt buộc với tài khoản quản trị và các thao tác nhạy cảm.

% =====================================================================
\section{Ứng dụng di động}
\label{sec:mobile}
% =====================================================================

Ứng dụng di động được phát triển bằng React~Native~\cite{reactnative2024} 0.76
kết hợp Expo~SDK~52~\cite{expo2024router}, viết hoàn toàn bằng
TypeScript~\cite{typescript_docs}. Lý do chọn React Native thay vì native
thuần là khả năng dùng chung phần lớn mã nguồn cho cả iOS và Android, giảm
đáng kể thời gian phát triển trong bối cảnh đồ án có giới hạn về nhân lực và
thời gian.

Expo Router~v4~\cite{expo2024router} được dùng làm hệ thống điều hướng dựa
trên cấu trúc thư mục (file-based routing), tương tự Next.js nhưng cho môi
trường native. Cấu trúc thư mục \texttt{app/} phản ánh trực tiếp cấu trúc
route: \texttt{(tabs)/} cho các màn hình tab chính, \texttt{(auth)/} cho luồng
đăng nhập, và các thư mục riêng cho allocation, maintenance, room detail và
violations.

Về quản lý trạng thái, ứng dụng phân biệt rõ server state và client state.
TanStack Query~v5~\cite{tanstack_query_docs} đảm nhiệm toàn bộ server state ---
fetch, cache, background refetch và optimistic updates. Zustand~v5~\cite{zustand_docs}
quản lý client state như thông tin người dùng đã xác thực và cấu hình giao
diện. HTTP communication dùng Axios~\cite{axios_docs} với interceptors tự động
gắn JWT access token vào header và xử lý refresh token khi nhận phản hồi 401.

Kết nối realtime từ ứng dụng di động đến backend được thiết lập qua
\texttt{socket.io-client}~\cite{socketio2024docs}, cho phép nhận thông báo phân
bổ và cập nhật trạng thái tức thì không cần polling. Token được lưu an toàn qua
\texttt{expo-secure-store}~\cite{expo_secure_store_docs}, dùng Keychain trên iOS
và Keystore trên Android. Mã QR thẻ cư trú được hiển thị bằng
\texttt{react-native-qrcode-svg}~\cite{rn_qrcode_docs}. Sentry React Native
được tích hợp để thu thập crash reports và performance traces từ phía di động.

Bên cạnh ứng dụng native, hệ thống còn triển khai thêm lớp Progressive Web App
ngay trên web portal, cho phép sinh viên cài đặt và dùng hệ thống như ứng dụng
độc lập mà không cần thông qua kho ứng dụng.

\subsection{Progressive Web App (PWA)}
\label{subsec:pwa}

Để mở rộng khả năng tiếp cận trên thiết bị di động mà không phụ thuộc hoàn
toàn vào ứng dụng native, hệ thống bổ sung lớp Progressive Web App
(PWA)~\cite{webdev_pwa} ngay trên web portal, dựa trên hai thành phần chuẩn
của web hiện đại: Service Worker và Web App Manifest~\cite{w3c_manifest}.

Service Worker (\texttt{public/sw.js}) đóng vai trò proxy mạng phía client,
áp dụng chiến lược cache phân biệt theo loại tài nguyên: cache-first cho asset
tĩnh (CSS, JavaScript, font, icon) để tăng tốc tải lại, và network-first cho
trang HTML động để đảm bảo dữ liệu luôn mới khi có kết nối. Khi mất mạng,
Service Worker phục vụ trang dự phòng \texttt{offline.html}; các yêu cầu nhạy
cảm như đăng nhập, API và Socket.IO được cấu hình bỏ qua cache hoàn toàn.

Web App Manifest (\texttt{public/manifest.webmanifest}) khai báo metadata để
hệ điều hành nhận diện web portal như một ứng dụng thực thụ: chế độ hiển thị
\texttt{standalone} ẩn thanh địa chỉ trình duyệt, bộ biểu tượng kèm icon
\texttt{maskable}, màu chủ đề, cùng các shortcut truy cập nhanh đến các màn
hình chính. Trên Android, hệ thống bắt sự kiện \texttt{beforeinstallprompt}
để hiển thị banner gợi ý cài đặt; trên iOS Safari, hệ thống hiển thị hướng
dẫn ``Thêm vào màn hình chính'' do nền tảng này chưa hỗ trợ cài đặt tự động.
Ngoài ra, một thanh trạng thái mạng phản hồi tức thì khi mất hoặc khôi phục
kết nối, và một thông báo xuất hiện khi có phiên bản Service Worker mới.

% =====================================================================
\section{Hạ tầng và vận hành}
\label{sec:infrastructure}
% =====================================================================

\subsection*{Container hóa}

Toàn bộ backend được đóng gói trong Docker~\cite{docker_docs} container dùng
multi-stage build với base image \texttt{node:20-alpine}. Stage đầu cài đặt
dependencies với \texttt{npm ci -{}-omit=dev}, stage hai sao chép source code
và chạy dưới non-root user \texttt{appuser} để tăng cường bảo mật. Các thư
mục không cần trong production như \texttt{mobile/}, \texttt{tests/} và
\texttt{scripts/} được loại bỏ để giảm kích thước image.

Docker Compose phục vụ triển khai local với hai service: \texttt{app} chạy
backend Node.js và \texttt{redis} chạy Redis~7-alpine. Service \texttt{app}
chỉ khởi động sau khi \texttt{redis} healthy. MongoDB Atlas được dùng như cơ
sở dữ liệu đám mây ngoài Docker Compose, kết nối qua biến môi trường
\texttt{MONGODB\_URI}.

\subsection*{Cache và truyền thông realtime}

Redis~\cite{redis2024pubsub} 5.x phục vụ hai mục đích: lưu trữ session phía
server thông qua \texttt{express-session} và cung cấp Socket.IO cluster adapter
(\texttt{@socket.io/redis-adapter}~\cite{socketio2024docs}), cho phép broadcast
events nhất quán qua nhiều instance backend. Khi Redis không khả dụng, hệ
thống fallback về chế độ local Socket.IO mà không crash. Endpoint
\texttt{GET /health} được triển khai để Docker và load balancer kiểm tra
trạng thái ứng dụng.

\subsection*{Email và thông báo ngoài ứng dụng}

Nodemailer~\cite{nodemailer_docs} 7.x gửi email qua SMTP (cấu hình bằng biến
\texttt{SMTP\_*}) phục vụ xác nhận đăng ký, thông báo kết quả phân bổ, gửi
OTP và hỗ trợ đặt lại mật khẩu. Trong triển khai hiện tại, SMTP được cấu hình
qua Gmail; cơ chế này có thể thay thế bằng dịch vụ SMTP doanh nghiệp khác mà
không cần sửa mã nguồn. Kênh thông báo chính trong ứng dụng vẫn là Socket.IO
realtime và hệ thống thông báo in-app, đảm bảo sinh viên nhận được thông tin
kịp thời ngay cả khi không kiểm tra email.

\subsection*{CI/CD}

GitHub Actions~\cite{github_actions} CI pipeline có hai job song song:
\texttt{backend} kiểm tra syntax toàn bộ file JavaScript trong \texttt{src/}
bằng \texttt{node -{}-check}; \texttt{mobile} chạy TypeScript type check bằng
\texttt{npx tsc -{}-noEmit}. Pipeline chạy tự động trên mọi push và pull
request, đảm bảo không có lỗi cú pháp hay type error nào lọt vào nhánh chính.
Trong môi trường production, PM2~\cite{pm2_docs} được dùng làm process manager
để giám sát và khởi động lại tự động khi có lỗi, với chế độ cluster tận dụng
nhiều nhân CPU.

% =====================================================================
\section{Công nghệ trực quan hóa 3D}
\label{sec:visualization}
% =====================================================================

Một trong những điểm khác biệt của hệ thống so với các giải pháp quản lý ký
túc xá hiện có là khả năng trực quan hóa không gian theo hai cấp độ: bản đồ
địa lý ba chiều ở cấp toàn khu và tham quan thực tế ảo ở cấp từng phòng. Hai
năng lực này được hiện thực bằng CesiumJS và Three.js, đều chạy trực tiếp trên
trình duyệt thông qua WebGL mà không cần cài plugin.

\subsection{CesiumJS --- Trực quan hóa địa lý 3D}
\label{subsec:cesium}

CesiumJS~\cite{cesiumjs} là thư viện JavaScript mã nguồn mở chuyên dựng địa
cầu và bản đồ ba chiều trên nền WebGL, được phát triển bởi Cesium~GS. Hệ thống
tích hợp CesiumJS vào tính năng ``Khám phá KTX 3D'' trên trang chủ, cho phép
người dùng quan sát vị trí thực tế của các khu ký túc xá trong không gian địa
lý. Mỗi tòa nhà được biểu diễn bằng marker tương tác; khi chọn marker, hệ
thống hiển thị thẻ thông tin kèm hình ảnh, mô tả và liên kết đến trang chi
tiết. Người dùng có thể phóng to, xoay và nghiêng góc camera mượt mà trên cả
máy tính lẫn thiết bị di động.

\subsection{Three.js --- Tour thực tế ảo phòng ký túc xá}
\label{subsec:threejs}

Three.js~\cite{threejs} là thư viện đồ họa 3D cấp cao trên nền WebGL, được
dùng để xây dựng tính năng tham quan phòng ký túc xá trước khi đăng ký. Mô
hình từng phòng được dựng ở định dạng glTF/GLB~\cite{khronos_gltf} --- tiêu
chuẩn mở của Khronos Group cho tài sản 3D trên web --- và nạp trực tiếp trong
trình duyệt qua \texttt{GLTFLoader}. Tính năng cung cấp hai chế độ: chế độ
tham quan có hướng dẫn di chuyển camera theo lộ trình định sẵn qua các điểm
quan trọng, và chế độ khám phá tự do cho phép người dùng tự điều khiển camera
bằng chuột hoặc cảm ứng kèm cơ chế xử lý va chạm cơ bản. Các mô hình được
nén bằng quantization và texture WEBP để giảm dung lượng truyền tải, phù hợp
với điều kiện kết nối di động thực tế.

Tính năng này giải quyết một hạn chế thực tế của các hệ thống ký túc xá hiện
tại: sinh viên thường phải đến tận nơi để xem phòng trước khi quyết định đăng
ký, đặc biệt bất tiện với sinh viên từ tỉnh xa. Tour 3D giúp sinh viên hình
dung diện tích, bố trí nội thất và điều kiện sinh hoạt ngay từ trang web, qua
đó giảm đáng kể nhu cầu khảo sát trực tiếp.

% =====================================================================
Chương 3 đã trình bày các lựa chọn công nghệ cùng lý do cụ thể đằng sau mỗi
quyết định: Node.js và MongoDB Atlas cho tính linh hoạt của mô hình dữ liệu và
hiệu năng I/O, Passport.js và JWT cho cơ chế xác thực đa lớp, React Native kết
hợp PWA cho khả năng tiếp cận trên nhiều thiết bị, Docker và GitHub Actions cho
quy trình vận hành nhất quán, và CesiumJS cùng Three.js cho trải nghiệm trực
quan hóa không gian. Các quyết định này tạo nền tảng cho chương 4 trình bày
thiết kế chi tiết, cách tổ chức mã nguồn và kết quả kiểm thử hệ thống.

\end{document}
